% C035 数学摘要；完整教学内容见正式六段正文。
$p>0$，$s=\pm1$，$A^2+B^2>0$。交点 $P(x_1,y_1),Q(x_2,y_2)$。
左右开口 $y^2=2spx$，优先用
\[
 Ay^2+2spBy+2spC=0.
\]
$A\ne0$ 时令 $H=p^2B^2-2spAC$，$\Delta_y=4H$；
$H>0$ 两不同实点，$H=0$ 相切（重计切点），$H<0$ 无实点。
\[
 y_1+y_2=-\frac{2spB}{A},\quad y_1y_2=\frac{2spC}{A};\qquad
 x_1+x_2=\frac{2spB^2-2AC}{A^2},\quad x_1x_2=\frac{C^2}{A^2}.
\]
横坐标固定关系为
\[
 A^2x^2+(2AC-2spB^2)x+C^2=0,\qquad \Delta_x=4B^2H.
\]
$B\ne0$ 时横坐标实根与交点一一对应；$B=0$ 时仅记录 $x=-C/A$，还须检验实纵坐标，不能仅凭横坐标重根判断相切。
$A=0$ 时，轴向直线恰有一个交点 $(sC^2/(2pB^2),-C/B)$，非相切。
有实弦时
\[
 S=\left(\frac{spB^2-AC}{A^2},-\frac{spB}{A}\right),\qquad
 PQ=\frac{2\sqrt{(A^2+B^2)H}}{A^2}.
\]
上下开口 $x^2=2spy$，优先用
\[
 Bx^2+2spAx+2spC=0.
\]
$B\ne0$ 时令 $K=p^2A^2-2spBC$，$\Delta_x=4K$；按 $K$ 正、零、负判断两点、相切、无实点。
\[
 x_1+x_2=-\frac{2spA}{B},\quad x_1x_2=\frac{2spC}{B};\qquad
 y_1+y_2=\frac{2spA^2-2BC}{B^2},\quad y_1y_2=\frac{C^2}{B^2}.
\]
纵坐标固定关系为 $B^2y^2+(2BC-2spA^2)y+C^2=0$；$A=0$ 时须另验实横坐标。
$B=0$ 时，轴向直线恰有一个交点 $(-C/A,sC^2/(2pA^2))$，非相切。
有实弦时
\[
 S=\left(-\frac{spA}{B},\frac{spA^2-BC}{B^2}\right),\qquad
 PQ=\frac{2\sqrt{(A^2+B^2)K}}{B^2}.
\]
$s=1$ 表示向右或向上，$s=-1$ 表示向左或向下。相切时长度公式给出零，不再是非零弦。
